% TOP_PROBLEM: {"id": 3260, "title": "Oriented trees in n-chromatic digraphs", "category_id": 3, "category": {"id": 3, "name": "graph_theory", "display_name": "Graph Theory", "description": "Problems involving graphs, networks, and their properties.", "slug": "graph-theory", "order_index": 3, "created_at": "2026-07-31T15:26:25.671Z"}, "status": "open", "problem_number": "OPG-46167", "set_id": 12}
% FIRST_POSTED: 2026-10-02
% ATTEMPT_STATUS: unresolved
% UnsolvedMath ID 3260; source catalogue code OPG-46167
% !TeX program = lualatex
% GPT-6 Astra Ultra research attempt; independently checked partial work, 2026-10-02.
\documentclass[11pt,a4paper]{article}
\usepackage[margin=25mm]{geometry}
\usepackage{fontspec}
\setmainfont{lmroman10-regular.otf}[BoldFont=lmroman10-bold.otf,
 ItalicFont=lmroman10-italic.otf,BoldItalicFont=lmroman10-bolditalic.otf]
\usepackage{amsmath,amssymb,mathtools}
\usepackage{unicode-math}
\setmathfont{latinmodern-math.otf}
\AtBeginDocument{\let\setminus\smallsetminus}
\usepackage{xurl,hyperref}
\hypersetup{colorlinks=true,urlcolor=blue}
\setcounter{secnumdepth}{0}
\setlength{\parindent}{0pt}
\setlength{\parskip}{5pt}
\setlength{\emergencystretch}{3em}
\begin{document}
\section*{Oriented trees in n-chromatic digraphs}\label{3260}
{\small UnsolvedMath ID 3260 · OPG-46167 · Graph Theory · OpenGarden}
\subsection{Definitions and mathematical statement}
Let $D=(V,A)$ be a finite loopless simple digraph, allowing opposite arcs
but not repeated arcs in one direction. Its underlying undirected graph
$U(D)$ has an edge $uv$ whenever at least one of $uv,vu$ is an arc.
The chromatic number used here is $\chi(D)=\chi(U(D))$: adjacent vertices
must have different colours, regardless of arc direction. It is not
the dichromatic number defined by acyclic colour classes.

An oriented tree $T$ is obtained by orienting each edge of a finite
undirected tree. Its order is $|V(T)|$. A copy of $T$ as a subdigraph of
$D$ means an injective map $\varphi:V(T)\to V(D)$ such that
$xy\in A(T)$ implies $\varphi(x)\varphi(y)\in A(D)$. Other arcs among
the image vertices are allowed, so this is not an induced-copy condition.

The conjecture asks whether, for every integer $k\ge2$, every digraph
$D$ with
\[
                            \chi(D)\ge2k-2
\]
contains every oriented tree of order $k$ in this sense. The quantifier
``every'' includes all tree shapes and all choices of directions.

The proposed threshold is sharp: a regular tournament on $2k-3$ vertices
has ordinary chromatic number $2k-3$ but every outdegree is $k-2$, so it
cannot contain the oriented star with one vertex pointing to $k-1$
leaves. No bound on indegrees, outdegrees, or the total order of $D$
is assumed beyond the chromatic hypothesis.
\subsection{Short English statement}
Does ordinary chromatic number at least 2k−2 force a digraph to contain every orientation of every k-vertex tree?
\subsection{Research attempt: all stars and a semidegree embedding lemma}
\paragraph{Scope and literature check.}
GPT-6 Astra Ultra, 2 October 2026. The chromatic number is that of the
underlying graph, and tree copies need not be induced. Bessy,
Gon\c calves and Reinald [S3] obtained a subquadratic sufficient
chromatic bound. Fan--Lu--Chen [S4], posted in September 2026, announce
a near-linear $O(k\log k)$ bound. This is a preprint result and still
differs from the exact conjectured $2k-2$ threshold. Here the work
establishes the threshold for every oriented star and identifies why
a straightforward greedy extension does not handle general trees.

\paragraph{1. Removing opposite arcs without losing chromatic number.}
Choose exactly one direction from each opposite pair, keeping all
other arcs. This produces an oriented graph $D'$ with precisely the
same underlying undirected graph as $D$. Every tree found in $D'$
also occurs in $D$. We may therefore assume in the proofs below that
incoming and outgoing neighbour sets are disjoint. This observation
is essential when the leaves of a mixed star must be distinct.

\paragraph{2. Stars with all arcs in one direction.}
A graph of chromatic number at least $2k-2$ has a vertex-critical
subgraph of chromatic number $2k-2$. Every vertex of that subgraph
has degree at least $2k-3$: otherwise a colouring of its deletion
with $2k-3$ colours extends to the vertex. Its orientation has at
least $(2k-3)|V|/2$ arcs. Average outdegree is consequently at least
$k-3/2$, and some integer outdegree is at least $k-1$. The vertex and
any $k-1$ of its outneighbours form the required outward star.
Applying the same argument to indegrees gives the inward star.
The critical subgraph is used only for counting; additional arcs
among leaves are harmless under the non-induced convention.

\paragraph{3. Every mixed star.}
Let the centre have $a\ge1$ outgoing leaves and $b\ge1$ incoming
leaves, where $a+b=k-1$. Suppose no vertex satisfies both
$d^+\ge a$ and $d^-\ge b$. Put
\[
 X=\{v:d^+(v)\le a-1\},\qquad Y=V\setminus X.
\]
Every vertex of $Y$ then has indegree at most $b-1$. For every
nonempty $W\subseteq X$, the induced underlying graph has at most
$(a-1)|W|$ edges, by counting outgoing arcs within $W$. Its average
degree is at most $2a-2$, so it has a vertex of degree at most that
number. Repeated deletion and greedy recolouring show
$\chi(D'[X])\le2a-1$. Counting incoming arcs gives, in exactly the
same way, $\chi(D'[Y])\le2b-1$.

Using disjoint palettes on the two sets yields
$\chi(D')\le2a+2b-2=2k-4$, a contradiction. Thus the required centre
exists. Its incoming and outgoing neighbour sets are disjoint in
$D'$, so selecting $a$ and $b$ leaves gives an injective star copy.
In fact mixed stars follow already from $\chi\ge2k-3$. Since every
tree on at most three vertices is a star, this also proves all
orientations in the cases $k=2,3$ of the original statement.

\paragraph{4. A sufficient extension condition and its obstruction.}
If a digraph has both minimum indegree and minimum outdegree at least
$k-1$, every oriented $k$-vertex tree embeds greedily. Root the tree
and process children after their parents. When embedding a child,
at most $k-2$ already occupied vertices other than the parent must
be avoided. The parent has at least $k-1$ neighbours in the required
direction, leaving an unused choice. No arcs between different
branches are prohibited.

Ordinary chromatic number does not furnish this semidegree condition,
even after passage to an induced subgraph. A transitive tournament
has arbitrarily large chromatic number, but every nonempty induced
subtournament has a source and a sink. It nevertheless contains every
oriented tree of suitable order, by placing a topological ordering
of the tree in its linear order. Thus the obstruction is to the
proposed extraction argument, not to the conjecture itself.

\paragraph{Verification, gap and final status.}
The script [S9] checks the regular-tournament star obstruction
below the proposed threshold and the source/sink obstruction just
used. The proofs cover all star sizes without computation. The first
remaining gap is coordinating extension at several internal tree
vertices under a chromatic hypothesis alone. No such invariant is
provided. Full target: UNRESOLVED; all oriented stars are verified.
\subsection{Sources}
\begin{itemize}
\item[S1] ULAM.AI and the UnsolvedMath Contributors, supplied problems.json, record 3260 (OPG-46167), OpenGarden collection. Catalogue identity and source transcription; CC BY 4.0.
\url{https://huggingface.co/datasets/ulamai/UnsolvedMath}.
\item[S2] Open Problem Garden, Oriented trees in n-chromatic digraphs. Original problem and discussion; the mathematical formulation below is expanded from the supplied source record.
\url{http://www.openproblemgarden.org/op/oriented_trees_in_n_chromatic_digraphs}.
\item[S3] St\'ephane Bessy, Daniel Gon\c calves and Amadeus Reinald, \emph{Oriented trees in $O(k\sqrt{k})$-chromatic digraphs, a subquadratic bound for Burr's conjecture}. \url{https://arxiv.org/abs/2402.19351}.
\item[S4] Liangdong Fan, Junying Lu and Yaojun Chen, \emph{A near-linear upper bound for Burr's conjecture} (preprint, 16 September 2026). \url{https://arxiv.org/abs/2609.18175}.
\item[S9] GPT-6 Astra Ultra, exact finite checks accompanying this attempt (2 October 2026). Repository script, Python standard library only. \texttt{scripts/check\_attempt\_batch03\_digraphs\_2026\_10\_02.py}.
\end{itemize}
\end{document}
